% !TeX root = ../main.tex
\begin{frame}{构型与逆向设计}
  \begin{SlideCanvas}
    \BodyAt{55}{84}{850}{支架结构设计从预设变形机制走向性能驱动的逆向设计，并连接单胞特性与器件指标。}
    \PicAt{55}{155}{323}{278}{assets/chiral2020_fig1.jpg}
    \CaptionAt{55}{444}{333}{Ruan等：手性单元与周向约束杆}
    \PointAt{425}{151}{475}{变形耦合构型}
    \DetailAt{425}{191}{475}{Wu 2018：反四手性与层级构型。\\Ruan 2020：周向约束与血管部署评价。}
    \draw[BICPageGray,line width=.75pt] (425,249)--(905,249);
    \PointAt{425}{270}{475}{多尺度拓扑优化}
    \DetailAt{425}{310}{475}{Xue 2020：共同考虑负泊松比与宏观刚度，完成成管分析及放大原型验证。}
    \draw[BICPageGray,line width=.75pt] (425,368)--(905,368);
    \PointAt{425}{389}{475}{单胞—器件性能映射}
    \DetailAt{425}{429}{475}{Carbonaro 2024：由等效刚度估计短缩率，并表征径向、轴向和扭转刚度。}
    \Sources{\RefWu；\RefRuan，图1。\\\RefXueA；\RefCarbonaro。前两项依据期刊摘要及原图。}
  \end{SlideCanvas}
  \BICConclusionBand{结构机制与性能映射共同建立了\BICBandEmph{面向器件需求的逆向设计基础}}
  \note{本页将原有构型史和逆向设计综述合并，围绕三项能力而非逐篇罗列。Wu与Ruan的描述限于公开摘要和期刊图件；Xue与Carbonaro已阅读全文。Xue原型放大15倍，聚合物材料用于相应机制验证。Carbonaro解析表征依赖线弹性及周期结构等假设。Clausen2015的有限变形泊松比方法另在课题组方法脉络中交代，完整18篇文献仍保留在引用说明中。}
\end{frame}
